After identifying potential investments, asset managers face the question about how to allocate the limited resources they have to each of them, such as to reduce risk while maintaining their target return. This task is of great importance for the asset managers' mandate to generate high risk adjusted returns. Harry M. Markowitz was the first to develop and present a systematic approach to this asset allocation problem in 1952, earning him a Nobel Memorial Prize in Economic Sciences in 1990. However, Markowitz's portfolio construction model presents some pitfalls in practical implementations that limit its performance out-of-sample. In this paper, we analyse Markowitz's Portfolio Construction and its pitfalls. Moreover, we present 2 approaches to improve asset allocations by tackling the main critique points of Markowitz's approach. Finally, we compare the results of the optimisation methods on the basis of multi-asset portfolios with the help of ETFs.